\documentclass[twoside]{article}

\usepackage{epsfig}
\usepackage{tabularx}
\usepackage{listings}
\usepackage{xcolor}
\usepackage{amsmath}

\definecolor{codegreen}{rgb}{0,0.6,0}
\definecolor{codegray}{rgb}{0.5,0.5,0.5}
\definecolor{codepurple}{rgb}{0.58,0,0.82}
\definecolor{backcolour}{rgb}{0.95,0.95,0.92}

\lstdefinestyle{mystyle}{
    backgroundcolor=\color{backcolour},   
    commentstyle=\color{codegreen},
    keywordstyle=\color{magenta},
    numberstyle=\tiny\color{codegray},
    stringstyle=\color{codepurple},
    basicstyle=\ttfamily\footnotesize,
    breakatwhitespace=false,         
    breaklines=true,                 
    captionpos=b,                    
    keepspaces=true,                 
    numbers=left,                    
    numbersep=5pt,                  
    showspaces=false,                
    showstringspaces=false,
    showtabs=false,                  
    tabsize=2
}

\lstset{style=mystyle}

\setlength{\oddsidemargin}{0.25 in}
\setlength{\evensidemargin}{-0.25 in}
\setlength{\topmargin}{-0.6 in}
\setlength{\textwidth}{6.5 in}
\setlength{\textheight}{8.5 in}
\setlength{\headsep}{0.75 in}
\setlength{\parindent}{0 in}
\setlength{\parskip}{0.1 in}

\newcommand{\lecture}[3]{
   \pagestyle{myheadings}
   \thispagestyle{plain}
   \newpage
   \setcounter{page}{1}
   \noindent
   \begin{center}
   \framebox{
      \vbox{\vspace{2mm}
    \hbox to 6.28in { {\bf CSE 291A/DSC 291:~Data Systems for Machine Learning, Spring 2026 \hfill} }
       \vspace{6mm}
       \hbox to 6.28in { {\Large \hfill #1  \hfill} }
       \vspace{6mm}
       \hbox to 6.28in{
        \it
        \begin{tabularx}{6.28in}{l>{\raggedleft\arraybackslash}X}
            Lecturer: #2 & Scribes: #3
        \end{tabularx}
       }
      \vspace{2mm}}
   }
   \end{center}
   \markboth{#1}{#1}
   \vspace*{4mm}
}

\begin{document}

\lecture{4: Tensor Formats, MatMul, Accelerators}{Hao Zhang}{
    Ian Jayachandran,
    Samar Karanch,
    Kevin Lin,
    Owen Park,
    Shaurya Raswan,
    Julie Wu,
    Jaewoong Yun
    } % Lecture name, Lecturer, Scribes

\section{Arithmetic Intensity (AI)}

Arithmetic Intensity is a measure of how compute-heavy an operation is relative to how much data needed to fetch from memory to perform that computation.

\textbf{Motivating Example:}

In standard multi-head attention, to get the Q, K, and V matrices, you need to compute each separately ($Q = X @ W_Q$, $K = X @ W_K$, $V= X @ W_V$).

Instead, you can speed up this process by performing a merged QKV by a single matrix operation ($X @ W_{QKV}$ and split into 3 matrices).

\subsection{Arithmetic Intensity Formula:}

\[
    \text{AI} = \text{\# operations } / \text{ \# bytes}
\]

This formula represents the ratio of compute work to memory work. A high AI value means that you are performing a lot of computations for each byte of memory read, while a low AI means the opposite. 

Transformer operations usually have low AI as you need to load entire matrices from memory while performing low-computation like dot products. This means the GPU will just sit idle waiting for loading data.

The practical goal is to raise AI enough so that memory is not a bottleneck in computing speed.

\subsection{Low AI Example 1:}
\begin{lstlisting}[language=C++, title={Simple add() kernel}]
void add(int n, float *A, float *B, float *C) {
    for (int i = 0; i < n; i++) {
        C[i] = A[i] + B[i];
    }
}
\end{lstlisting}

The code computes the operation of $A[i] + B[i]$ $n$ times via the for loop. 

It also loads $A[i], B[i]$ from memory $n$ times each and stores $C[i]$ to memory $n$ times each, so $3n$ computations total.

According to the AI formula, we have $\#$ operations as $n$, and $\#$ memory loads as $3n$. Therefore, $AI = \frac{n}{3n} = \frac{1}{3}$.


\subsection{Low AI Example 2:}

\begin{lstlisting}[language=C++, title={Simple add() and mul() kernel}]
void add(int n, float* A, float* B, float* C){
    for (int i=0; i<n; i++) {
        C[i] = A[i] + B[i];
    }
}

void mul(int n, float* A, float* B, float* C) {
    for (int i=0; i<n; i++) {
        C[i] = A[i] * B[i];
    }
}

float* A, *B, *C, *D, *E, *tmp1, *tmp2;
// assume arrays are allocated here
// compute E = D + ((A + B) * C)
add(n, A, B, tmp1);
mul(n, tmp1, C, tmp2);
add(n, tmp2, D, E);
\end{lstlisting}

Both the add and multiply operations are 1/3 as explained in Example 1. So, if we call the add operation twice and multiply operation once, we have $3n$ total computations and $9n$ total loads and stores. So, $\frac{2n}{6n} = \frac{1}{3}$

\textbf{Therefore, AI is still $\frac{1}{3}$.} 

This is wasteful as we see $tmp1$ being written to memory during add(n, A, B, tmp1), and loaded from memory during mul(n, tmp1, C, tmp2). The same process applies to $tmp2$. Why not keep $tmp1$ or $tmp2$ in memory for quicker processing?

\subsection{More Efficient Example (Kernel Fusion):}


\begin{lstlisting}[language=C++, title={Fused kernel}]
float *A, *B, *C, *D, *E, *tmp1, *tmp2;
// assume arrays are allocated here
// compute E = D + ((A + B) * C)
void fused(int n, float *A, float B, floatC, float *D, float *E) {
    for (int i = 0; i < n; i++) {
        E[i] = D[i] + (A[i] + B[i]) * C[i];
    }
}
// compute E = D + ((A + B) * C)
fused(n, A, B, C, D, E);
\end{lstlisting}

We can instead combine all compute operations in one kernel as shown above. 

Now we have the following:

$\#$ $Operations = 3n$ (2 additions, 1 multiplication for each loop iteration) 

$\#$ $Bytes$ $Loaded = 5n$ (For each iteration, load $A[i], B[i], C[i], D[i]$, and store once to $E[i]$) 

\textbf{AI = 3/5}. Higher AI! Since we lowered the amount of memory accesses, we have sped up our kernel processing time.

\section{Graph Optimization}
\subsection{Parallelization}
\subsection{Parallelization Problems}
- partition
- communicate
- schedule
- consistency
- auto parallelize
\subsection{Runtime \& Scheduling}
- schedule compute / communication / memory as fast as possible
- overlap communication w/ compute
- subject to memory constraints
\subsection{Operator Implementation}
- get fastest possible implementation of MatMul, Conv2d
- optimized for diff GPU types: V100, A100, H100, phone, TPU
deploy model for diff GPU types need adaptation based on varying constraints
- diff precision fp32, fp16, fp8, fp4
- diff shape conv2d\_3x3, conv2d\_5x5, matmul2D, 3D, attention
- goal is to maximize AI

\section{Fast Operators}
\subsection{Vectorization}


\begin{lstlisting}[language=C++, title={Unvectorized}]
float A[256], B[256], C[256]
for (int i = 0; i < 256; ++i) {
    C[i] = A[i] + B[i]
}
\end{lstlisting}


\begin{lstlisting}[language=C++, title={Vectorized}]
for (int i = 0; i < 64; ++i) {
    float4 a = load_float4(A + 1*4);
    float4 b = load_float4(B + i*4);
    float4 c = add_float4(a, b);
    store_float4(C + i* 4, c);
\end{lstlisting}

Vectorization minimizes writes

\subsection{Data Layout}
- store matrix in memory
- stored sequentially (1D, no tensor awareness)
- computation memory limit crucial for quantized ML
- Row Major: A[i, j] = A.data[i * A.shape[1] + j] (left to right)
- Col Major: A[i, j] = A.data[j * A.shape[0] + i] (top to bottom)
- pytorch utilizes row major calc

INCLUDE EXAMPLE PROGRAM ROW/COL MAJOR SETUP
setup row major 2D array, but reading in column major, swap to make program efficient

\subsubsection{Strides}
ML systems store data in strides format: A[i, j] = A.data[offset + i * A.strides[0] + j * A.strides[1]]

offset: offset of tensor relative to underlying storage
strides: strides[i] = \# of elements needed to be skipped in memory to move "one unit" in the i-th dimeniosn of the tensor

\begin{lstlisting}[language=Python]

\end{lstlisting}

A.strides[0] = 1, A.strides[1] = A.shape[0] --> col major
A.strides[0] = A.shape[1], A.strides[1] = 1 --> row major

if a tensor of shape [1,2,3,4] is stored in contiguous in memory following row major, its strides are (24, 12, 4, 1)

strides can separate underlying storage \& view of tensor
enable zero-copy of frequently used operators

\subsubsection{Operator: Slice}
read first 3 rows, 2 columns
change offset by +1, reduce shape to [3,2]
no copy needed of tensor
\subsubsection{Operator: Transpose}
permute by moving first stride to end 
zero copy, underlying storage unchanged
\subsubsection{Operator: Broadcast}
b.strides=[1], b.shape=[1, 3], b.data=[1, 2, 3]
broadcast so b.shape=[4, 3], b.data=[1, 2, 3], b.strides=[?]
add elements to stride tensor
\subsubsection{Operator: Swap Tiles}

\subsubsection{Problems of Strides}
- memory access can become not continuous
- many vectorized ops require continuous storage

\subsection{Parallelization}
elementise, deeper dive will be covered in later lecture
ADD EXAMPLE CODE FROM LECTURE
parallelize through threads
on gpu stream through processors

\section{MatMul}
\begin{lstlisting}[language=C++]
float A[n][n], B[n][n], C[n][n];

for (int i = 0; i < n; ++i) {
    for (int j = 0; j < n; ++j) {
        C[i][j] = 0;
        for (int k = 0; k < n; ++k) {
            C[i][j] += A[i][k] * B[j][k];
        }
    }
}
\end{lstlisting}

\subsection{MatMul Complexity}
O(N\^3) complexity, but can get O(N\^2.371552)
make matmul fast by increasing ops, decreasing bytes

ideally we want everythign to be local to processors (in registers), but registers are expensive \& small, thus we have memory hierarchy (L1, L2, etc. DRAM) on increasing latency

\subsection{Register Tiled MatMul}
\begin{lstlisting}[language=C++]
dram float A[n][n], B[n][n], C[n][n];

for (int i = 0; i < n; ++i) {
    for (int j = 0; j < n; ++j) {
        register float c = 0;
        for (int k = 0; k < n; ++k) {
            register float a = A[i][k];
            register float b = B[j][k];
            c += a * b;
        }
        C[i][j] = c;
    }
}
\end{lstlisting}
read a = n\^3, read b = n\^3, write C = n\^2
need 3 registers (1 a, 1 b, 1 c) = 1 + 1 + 1 = 3
total read cost = 2*n\^3*speed(moving component from dram to register)
2x bc for matmul you multiply a*b (1 op) and then add to c (1 op) = 2 ops

\begin{lstlisting}[language=C++]
dram float A[n/v1][n/v3][v1][v3];
dram float B[n/v2][n/v3][v2][v3];
dram float C[n/v1][n/v2][v1][v2];

for (int i = 0; i < n/v1; ++i) {
    for (int j = 0; j < n/v2; ++j) {
        register float c[v1][v2] = 0;
        for (int k = 0; k < n/v3; ++k) {
            register float a[v1][v3] = A[i][k];
            register float b[v2][v3] = B[j][k];
            c += dot(a, b.T);
        }
        C[i][j] = c;
    }
}
\end{lstlisting}
read a = n\^3/v2, read b = n\^3/v1, write C = n\^2
need v1*v3 + v2*v3 + v1*v2 registers
total read cost = (n\^3/v2 + n\^3/v1)*speed(dram -> register)
\subsubsection{Cache-aware Tiling}
\begin{lstlisting}[language=C++]
dram float A[n/b1][b1][n];
dram float B[n/b2][b2][n];
dram float C[n/b1][n/b2][b1][b2];

for (int i = 0; i < n/b1; ++i) {
    l1cache float a[b1][n] = A[i];
    for (int j = 0; j < n/b2; ++j) {
        l1cache b[b2][n] = B[j];
        
        C[i][j] = dot(a, b.T);
    }
}
\end{lstlisting}
data movement path: DRAM -> L1 Cache (cache-tiling) -> register (register tiling)
A dram -> l1 cost: n / b1 * n * b1 = n\^2
B dram -> l1 cost: n / b1 * n / b2 * 2 * n = n\^3 / 1
b1 * n + b2 * n < L1 cache size 

\begin{lstlisting}[language=C++]
dram float A[n/b1][b1/v1][n][v1];
dram float B[n/b2][b2/v2][n][v2];

for (int i = 0; i < n/b1; ++i) {
    l1cache float a[b1/v1][n][v1] = A[i];
    for (int j = 0; j < n/b2; ++j) {
        l1cache b[b2/v2][n][v2] = B[j];
        for (int x = 0; x < b1/v1; ++x) {
            for (int y = 0; y < b2/v2; ++y) {
                register float c[v1][v2] = 0;
                for (int k = 0; k < n; ++k) {
                    register float ar[v1] = a[x][k][:];
                    register float br[v2] = b[y][k][:];
                    C += dot(ar, br.T)
                }
            }
        }
    }
}
\end{lstlisting}
outside loop = cache-tiling
- use b1, b2
- DRAM -> L1 reads here
- NEED COST EQUATIONS
inside loop = register tiling
- use v1, v2
- L1 -> register cache reads here
- NEED COST EQUATIONS

\begin{lstlisting}[language=C++]
float A[n][n];
float B[n][n];
float C[n][n];

C[i][j] = sum(A[i][k] * B[j][k], axis=k)
\end{lstlisting}

\subsection{Homework Candidate}

\end{document}